\section{El fondo cosmológico en UG con difusión}
\label{sec:fondo}

Con las ecuaciones de la sección anterior se puede responder la primera pregunta del trabajo: cómo entra $Q$ en la cosmología de fondo. La derivación es corta, y su resultado es simple e importante para todo lo que sigue.

\subsection{Las ecuaciones de Friedmann con $Q$}

Se toma la métrica de FLRW en tiempo cósmico \eqref{eq:FLRW} y se escribe la ecuación de campo \eqref{eq:EinsteinUG} en forma mixta, $G^a{}_b+\lambda\,\delta^a_b=\kap\,T^a{}_b$. En el fondo homogéneo, $\lambda$ depende solo del tiempo, $\Lb(t)=\Lambda_0+\kap\,Q(t)$. Con $G^a{}_b$ de \eqref{eq:Gfondo} y $T^0{}_0=-\rho$, $T^i{}_j=p\,\delta^i_j$, las componentes $00$ e $ij$ dan \etq{P}
\begin{equation}
3H^2=\kap\rho+\Lb,\qquad 2\dot H+3H^2=-\kap p+\Lb,
\label{eq:FriedUG}
\end{equation}
que coinciden con las de \cite{leon22}. Si la materia es un inflatón, con $\rho$ y $p$ de \eqref{eq:rhopphi}, resultan
\begin{align}
3\MP^2H^2&=\tfrac12\dot\phi^2+V+Q+\Lambda_0\MP^2, \label{eq:H2UG}\\
\dot H&=-\frac{\dot\phi^2}{2\MP^2}. \label{eq:HdotUG}
\end{align}
La \eqref{eq:HdotUG} sale de restar las dos ecuaciones de \eqref{eq:FriedUG}, y es \emph{idéntica a la de RG}, \eqref{eq:HdotPhi}: la difusión $Q$ no aparece en $\dot H$. La \eqref{eq:H2UG}, en cambio, tiene dos términos nuevos en el lado derecho: $Q$ y la constante $\Lambda_0$. Esta asimetría entre las dos ecuaciones es el resultado central de la sección, y su explicación es simple. El término $Q$ se comporta como un fluido con densidad $\rho_Q=Q$ y presión $p_Q=-Q$ (ecuación de estado $w=-1$), y como $\dot H=-\tfrac{\kap}{2}(\rho+p)$ solo ve la suma $\rho+p$, un fluido con $\rho_Q+p_Q=0$ no puede aparecer en $\dot H$ aunque sí en $H^2$.

La continuidad se obtiene derivando la primera ecuación de \eqref{eq:FriedUG} y usando la segunda:
\begin{equation}
\dot\rho+3H(\rho+p)=-\dot Q.
\label{eq:contUG}
\end{equation}
Es la componente temporal de $\nabla_aT^{ab}=\nabla^bQ$: la materia ordinaria \emph{no} conserva su energía, y lo que pierde (o gana) lo intercambia con el sector de $Q$. La energía total $\rho+Q$ sí se conserva. Si $Q$ decrece ($\dot Q<0$, que es el caso en el modelo numérico), la materia gana energía.

\subsection{La constante de integración}

La constante $\Lambda_0$ no es un parámetro de la acción sino un dato inicial. Evaluando \eqref{eq:H2UG} en un instante inicial, se obtiene \etq{P}
\begin{equation}
\Lambda_0=3H_{\rm ini}^2-\frac{\rho_{\rm ini}+Q_{\rm ini}}{\MP^2},
\label{eq:Lambda0}
\end{equation}
es decir, la constante que asegura que la ecuación de vínculo de Friedmann se cumpla en el estado inicial elegido. Dos observaciones. Primero, esta es una diferencia sustancial con RG: en RG, si se quiere comparar dos universos con el mismo potencial y las mismas condiciones $(\phi,\dot\phi)$ iniciales, $H_{\rm ini}$ queda fijado por la ecuación de Friedmann; en UG, $H_{\rm ini}$ y $\Lambda_0$ son cantidades relacionadas y se puede elegir una u otra libremente. Segundo, esa libertad hace que la pregunta ``¿qué es el mismo modelo en UG y en RG?'' tenga más de una respuesta razonable, y se vuelve una decisión que condiciona la magnitud de los efectos (sección~\ref{sec:decisiones}).

\subsection{La ecuación del inflatón}
\label{sec:ecinflaton}

Falta cerrar el sistema, porque \eqref{eq:H2UG} y \eqref{eq:HdotUG} hay que complementarlas con la ecuación de movimiento del campo. Como $\dot\rho_\phi+3H(\rho_\phi+p_\phi)=\dot\phi\,(\ddot\phi+3H\dot\phi+V')$, si se atribuye \emph{toda} la no conservación de \eqref{eq:contUG} al inflatón se obtiene \etq{P, S}
\begin{equation}
\ddot\phi+3H\dot\phi+V'(\phi)=-\frac{\dot Q}{\dot\phi}.
\label{eq:KGQ}
\end{equation}
Con $Q$ constante se recupera la ecuación de Klein--Gordon de RG. Esta ecuación no se deduce de la acción \eqref{eq:Sphi}: como se explicó en el recuadro de la sección~\ref{sec:difusionQ}, un inflatón con acción covariante estándar conserva su energía. Es una ecuación fenomenológica que supone que $Q$ intercambia energía \emph{solo con el inflatón}. La alternativa que usa \cite{piccirilli23} es distinta: el inflatón se conserva y $Q$ alimenta un fluido de radiación. Con un solo campo, \eqref{eq:KGQ} es la única manera de que \eqref{eq:H2UG}, \eqref{eq:HdotUG} y \eqref{eq:contUG} sean consistentes, porque si $\phi$ obedeciera la ecuación de Klein--Gordon estándar mientras $Q$ varía, la derivada de \eqref{eq:H2UG} no coincidiría con \eqref{eq:HdotUG}. Por eso es la que se implementa en el cálculo numérico. Nada de lo que se deriva sobre los tensores en la sección~\ref{sec:tensores} depende de esta elección, porque en el sector tensorial solo entra la suma de todos los contenidos.

\subsection{El primer parámetro de flujo y el slow-roll con $Q$}

Con \eqref{eq:H2UG} y \eqref{eq:HdotUG}, el primer parámetro de flujo $\epsilon_1=-\dot H/H^2=\dot\phi^2/(2\MP^2H^2)$ es \etq{P}
\begin{equation}
\epsilon_1^{\rm UG}=\frac{3\dot\phi^2}{\dot\phi^2+2\,(V+Q+\Lambda_0\MP^2)},\qquad
\epsilon_1^{\rm RG}=\frac{3\dot\phi^2}{\dot\phi^2+2V}.
\label{eq:eps1UG}
\end{equation}
La aceleración ($\epsilon_1<1$) equivale a $\dot\phi^2<V+Q+\Lambda_0\MP^2$; para $\Lambda_0=0$ es la condición ``$Q+V>2X$'' con $X=\dot\phi^2/2$ que da \cite{piccirilli23}. Para un fluido perfecto con $p=w\rho$ y $\Lambda_0=0$, definiendo $\Gamma\equiv Q/\rho$, se obtiene $\epsilon_1=3(1+w)/[2(1+\Gamma)]$, que reproduce el caso de radiación de \cite{leon22}, $\epsilon_1=2/(1+\Gamma)$.

Esta aproximación usa Klein--Gordon conservativa: para el único escalar que recibe la transferencia exige $\dot Q=0$. Con una ley local fija, $U=V+Q(\phi)+\Lambda_0 M_P^2$, la aproximación correspondiente usa $U$ y $U\prime$. El código integra la ecuación completa con fuente para $Q(N)$.

En el régimen de slow-roll, con $|\ddot\phi|\ll|3H\dot\phi|$ y la ecuación de Klein--Gordon estándar \etq{P, S}:
\begin{equation}
H^2\simeq\frac{V+Q+\Lambda_0\MP^2}{3\MP^2},\qquad
\epsilon_1\simeq\frac{\MP^2}{2}\Big(\frac{V'}{V+Q+\Lambda_0\MP^2}\Big)^2.
\label{eq:slowUG}
\end{equation}
Si $Q$ es constante, es exactamente la RG con el potencial desplazado en una constante, $V\to V+\text{cte}$; es coherente con la afirmación de \cite{ellis14} de que, con las ecuaciones sin traza, ``la inflación procede como de costumbre'' (lo que solo se leyó en el resumen). Todo el efecto de una $Q$ variable está en que $H^2$ y $\epsilon_1$ cambian respecto de RG con el mismo potencial, y ese será el canal por el que, como se verá, $Q$ llega al espectro tensorial.

\begin{table}[t]
\centering\small
\renewcommand{\arraystretch}{1.3}
\begin{tabular}{@{}p{3.1cm}p{3.5cm}p{3.7cm}p{4.2cm}@{}}
\toprule
& \textbf{RG} & \textbf{UG conservativa} ($Q$ cte) & \textbf{UG no conservativa}\\
\midrule
$3\MP^2H^2$ & $\rho_\phi\;(+\Lambda\MP^2)$ & $\rho_\phi+\Lambda_0\MP^2$ & $\rho_\phi+Q(t)+\Lambda_0\MP^2$\\
$\dot H$ & $-\dot\phi^2/2\MP^2$ & igual & \textbf{igual}\\
Origen de la constante & parámetro de la acción & integración, \eqref{eq:Lambda0} & integración; $Q(t)$ libre\\
Ecuación del inflatón & Klein--Gordon & Klein--Gordon & \eqref{eq:KGQ} si $Q$ se acopla al inflatón\\
Equivalencia con RG & --- & \textbf{exacta}, $\Lambda=\Lambda_0$ & no: $H(t)$ distinto para el mismo $V$\\
\bottomrule
\end{tabular}
\caption{Comparación de las ecuaciones de fondo con un inflatón.}
\label{tab:fondo}
\end{table}

La tabla~\ref{tab:fondo} resume la comparación. La conclusión de esta sección para la primera pregunta del trabajo es que el vínculo unimodular modifica el fondo cuando hay difusión, a través de $H^2$ y por lo tanto de $\epsilon_1$, pero no modifica la relación $\dot H=-\dot\phi^2/2\MP^2$; y que en el caso conservativo la modificación se reduce a una constante cosmológica de integración. Las derivaciones existentes de este fondo en la literatura \cite{leon22,piccirilli23} suponen radiación como contenido de materia; el resultado de esta sección lo obtiene para un inflatón y sin suponer nada sobre $Q$ salvo que es un escalar.
